%---------------------- TEMPLATE FOR REPORT ------------------------------------------------------------------------------------------------------%

%\thispagestyle{fancy}   % 插入页眉页脚                                        %
%%%%%%%%%%%%%%%%%%%%%%%%%%%%% 用 authblk 包 支持作者和E-mail %%%%%%%%%%%%%%%%%%%%%%%%%%%%%%%%%
%\title{More than one Author with different Affiliations}				     %
%\title{\rm{VASP}的电荷密度存储文件\rm{CHGCAR}}
%\title{面向高温合金材料设计的计算模拟软件中的几个主要问题}
%\title{计算材料团队本年度工作规划:~\\经营计划与科研计划}
\title{平生风义兼师友}
\author[ ]{}   %
%\author[ ]{姜~骏\thanks{jiangjun@bcc.ac.cn}}   %
%\affil[ ]{北京市计算中心}    %
%\author[a]{Author A}									     %
%\author[a]{Author B}									     %
%\author[a]{Author C \thanks{Corresponding author: email@mail.com}}			     %
%%\author[a]{Author/通讯作者 C \thanks{Corresponding author: cores-email@mail.com}}     	     %
%\author[b]{Author D}									     %
%\author[b]{Author/作者 D}								     %
%\author[b]{Author E}									     %
%\affil[a]{Department of Computer Science, \LaTeX\ University}				     %
%\affil[b]{Department of Mechanical Engineering, \LaTeX\ University}			     %
%\affil[b]{作者单位-2}			    						     %
											     %
%%% 使用 \thanks 定义通讯作者								     %
											     %
\renewcommand*{\Authfont}{\small\rm} % 修改作者的字体与大小				     %
\renewcommand*{\Affilfont}{\small\it} % 修改机构名称的字体与大小			     %
\renewcommand\Authands{ and } % 去掉 and 前的逗号					     %
\renewcommand\Authands{ , } % 将 and 换成逗号					     %
\date{} % 去掉日期									     %
%\date{2020-12-30}									     %

%%%%%%%%%%%%%%%%%%%%%%%%%%%%%%%%%%%%%%%%%%  不使用 authblk 包制作标题  %%%%%%%%%%%%%%%%%%%%%%%%%%%%%%%%%%%%%%%%%%%%%%
%-------------------------------The Title of The Report-----------------------------------------%
%\title{报告标题:~}   %
%-----------------------------------------------------------------------------

%----------------------The Authors and the address of The Paper--------------------------------%
%\author{
%\small
%Author1, Author2, Author3\footnote{Communication author's E-mail} \\    %Authors' Names	       %
%\small
%(The Address，City Post code)						%Address	       %
%}
%\affil[$\dagger$]{清华大学~材料加工研究所~A213}
%\affil{清华大学~材料加工研究所~A213}
%\date{}					%if necessary					       %
%----------------------------------------------------------------------------------------------%
%%%%%%%%%%%%%%%%%%%%%%%%%%%%%%%%%%%%%%%%%%%%%%%%%%%%%%%%%%%%%%%%%%%%%%%%%%%%%%%%%%%%%%%%%%%%%%%%%%%%%%%%%%%%%%%%%%%%%
\maketitle
%\thispagestyle{fancy}   % 首页插入页眉页脚 
2026年3月3号是农历的正月十五，这天的天气很阴沉，北京城区还飘了一点雪花，后来我才知道，就是在这天中午王崇愚先生离开了我们。我获悉王先生离世的消息已经是第二天的晚上，是王殿武老师告知我的。2025年9月底，曾有过一次到钢研总院讨论的机会，那次我曾询问过殿武老师王先生的近况，殿武老师说:``王老师在\textrm{ICU}。''我便没有再深问。从那以后，我告诉自己，没有王先生的消息，就是好消息。不想噩耗却还是来的这么快。

最初拜识王先生是2012年初秋，当时我在一个叫北京宏剑公司的科学计算模拟软件代理公司上班，公司聘请了不少计算化学、计算材料学的老院士做顾问，每逢节假日前，都会礼节性地拜望老先生们，借此联络感情。那次记得是清华刚开学不久，我们在王先生清华物理系三楼办公室见面。因为是初次面谒，只是简单寒暄，但王先生听说我读过一些计算软件的代码后，表现出了很大的兴趣，还问了我的师承。从先生办公室离开的时候，我根本没有想到，在此后十多年里，竟能得到王先生如此多的帮助和提携。幸运的是，第二次见王先生就有机会和先生作了比较深入的交谈。同年的11月，宏剑公司在天津大学举办了一次\textrm{VASP}软件培训班，王先生也参加了。在去天津的路上，我作为陪同人员，和王先生说起之前学习计算材料方法的经历，王先生知道我是学\textrm{LAPW}方法出身的，现在开始接触\textrm{VASP}软件和\textrm{PAW}方法，就建议我说，``你之前能通过读代码的方式学\textrm{WIEN2k}软件，那\textrm{VASP}最好也能用这个办法拿起来，国内真正懂方法、懂代码的人不多。''，除了具体的材料计算方法，王先生也表达了对相关科学计算软件的国产化问题的关切，谈了很多自己的看法，记得先生还着意强调了``\textrm{innovation}''和``creation''两者词义上的区别。一路上王先生的谈兴很浓，看得出来，怹对软件和方法是很有的想法。记得我那次培训，我的任务主要介绍的是\textrm{VASP}软件的\textrm{PAW}方法和代码的总体结构。报告结束之后，王先生又专门把我叫到身边，再次关照我，务必花一点精力，把\textrm{VASP}的代码``捋''一遍。老先生的信任与鼓励，让我觉得既荣幸，又惶恐，怹大约看出了我的犹豫，又说:``以后有问题，可以找我讨论。你是黎乐民的学生，黎老师我是了解的，我有软件方面的问题也可以找你''，然后还和我互留了手机号。更让我没有想到的是，王先生以极大的热忱，随同我们公司成员全程参与了这次两天的软件培训。一位科学院资深院士，和很多研究生一起，参与一个计算软件的培训课程，并且认真记录，多次提问，这样的场景，今天想来都很让人感动。回京路上，我和王先生明显亲近了许多，话题也从科学计算渐渐发散开去，我无意中说起和北大技物系王德民教授除了师生关系外，还有点``忘年交''的意思，王先生更高兴了，说怹和王德民老师最要好了。可是说起来二位先生虽然同住蓝旗营8号楼，因为在不同单元，已经是多年未见了。先生无意中的

春节的时候，借着给师母微信拜年的时候，也问候了王先生，
